\begin{table}[htbp!]
\centering
\caption{F1-\textit{score} por categoria de consulta e por modelo}
\label{tab:resultados_por_categoria}
\small
\begin{tabular}{|l|c|c|}
\hline
\textbf{Categoria} & \textbf{Qwen 3 4B} & \textbf{Gemma 4 E2B} \\
\hline
Simples (n=12) & 0.889 & 1.000 \\
Compostas (n=45) & 0.692 & 0.696 \\
Com código MI/INOM (n=16) & 0.863 & 0.815 \\
Com referência temporal relativa (n=16) & 0.669 & 0.489 \\
Com ordenação (n=11) & 0.723 & 0.633 \\
Ambíguas / variações ortográficas (n=38) & 0.649 & 0.664 \\
\hline
\textbf{Média ponderada} & 0.706 & 0.672 \\
\hline
\end{tabular}
\fonte{Elaborado pelo autor. Uma consulta pode pertencer a mais de uma categoria; n = consultas da categoria nas métricas principais.}
\end{table}
